% Choose environments by purpose; the theme supplies their appearance.
\begin{frame}[fragile]{Everyday blocks}
  \framesubtitle{Give a result, highlight a limitation, or show an example.}
  \small
  \begin{columns}[T,onlytextwidth]
    \begin{column}{.47\textwidth}
      \begin{block}{Main result}
        State the finding in one sentence.
      \end{block}
      \begin{alertblock}{Limitation}
        Explain what affects the interpretation.
      \end{alertblock}
      \begin{exampleblock}{Example}
        Show a concrete application.
      \end{exampleblock}
    \end{column}
    \begin{column}{.49\textwidth}
      {\footnotesize
\begin{verbatim}
\begin{block}{Main result}
  State the finding.
\end{block}

\begin{alertblock}{Limitation}
  Explain the limitation.
\end{alertblock}

\begin{exampleblock}{Example}
  Show an application.
\end{exampleblock}
\end{verbatim}
      }
      \smallskip
      Titles use braces; put your content inside the environment.
    \end{column}
  \end{columns}
\end{frame}

\begin{frame}[fragile]{Context, remarks, and quotations}
  \framesubtitle{Amurmaple adds information and remark boxes.}
  \small
  \begin{columns}[T,onlytextwidth]
    \begin{column}{.47\textwidth}
      \begin{information}[Data quality]
        This panel gives context needed to understand the result.
      \end{information}
      \begin{remark}[A useful caveat]
        Add a brief observation without interrupting the main argument.
      \end{remark}
      \smallskip
      Square brackets add an optional title or note. Omit them for the
      default heading.
    \end{column}
    \begin{column}{.49\textwidth}
      {\footnotesize
\begin{verbatim}
\begin{information}[Data quality]
  Add useful context.
\end{information}

\begin{remark}[A useful caveat]
  Add a brief observation.
\end{remark}
\end{verbatim}
      }
      \smallskip
      \textbf{Other useful formats}\par
      {\footnotesize
\begin{verbatim}
\begin{quotation}[Source]
  A short quotation.
\end{quotation}
\begin{abstract}
  A brief overview.
\end{abstract}
\end{verbatim}
      }
    \end{column}
  \end{columns}
\end{frame}

\begin{frame}[fragile]{Mathematical statements}
  \framesubtitle{Use the standard Beamer environments for structured arguments.}
  \small
  \begin{columns}[T,onlytextwidth]
    \begin{column}{.47\textwidth}
      \begin{definition}[Sampling interval]
        \(\Delta t\) is the time between samples.
      \end{definition}
      \begin{theorem}[Positive sample rate]
        If \(\Delta t>0\), then \(1/\Delta t>0\).
      \end{theorem}
      \begin{proof}
        The reciprocal of a positive number is positive.
      \end{proof}
    \end{column}
    \begin{column}{.49\textwidth}
      {\footnotesize
\begin{verbatim}
\begin{theorem}[Positive rate]
  State the claim.
\end{theorem}
\begin{proof}
  Give the argument.
\end{proof}
\end{verbatim}
      }
      \smallskip
      \textbf{Use the same pattern for}\par
      \texttt{definition}, \texttt{lemma}, \texttt{corollary},
      \texttt{fact}, \texttt{problem}, and \texttt{solution}.
      \medskip\par
      The bracketed description is optional. Replace both environment
      names when adapting an example.
    \end{column}
  \end{columns}
\end{frame}
